% Transcription — fonds Alexandre Grothendieck, shelfmark 19, batch 3 (pages 41-60).
% Established from the facsimile archives/batches/19/batch-03.pdf (a mirror of
% https://grothendieck.umontpellier.fr/19.pdf), per the skill
% transcribe-grothendieck.
%
% Pages 49, 51 and 55 are blank sheets — absent.
%
% Three pieces. Pages 41-42 close the "Remarques — SGA 4" set begun in batch 2,
% but fall outside the "11 p" the page 28 wrapper counts (pages 30-40): page 41
% is a faint pencil remark headed XII on how little the quasi-compactness and
% quasi-separation hypotheses are really used in 6.5-6.7, and page 42 is a
% typescript sheet of the exposé whose part (iii) is struck through, carrying
% at its foot a handwritten instruction for redoing the proof. Page 41 is by
% some way the hardest sheet of the batch and is transcribed with many gaps.
%
% Page 43 is the wrapper of the second piece and names it: "Foncteurs
% localisant et catégories de fractions", with a pencilled tally "4 p + 5 p".
% Pages 44-53 are the theory of quotient categories, written as numbered items
% and not in reading order — page 44 carries items 3) and 4), page 45 items 2)
% and 1). What is worked out is the section functor S adjoint to the quotient
% T : A -> A/C, the C-closed objects and the C-isomorphisms, when the
% inclusion carries injectives to injectives, and (pages 46-47) the two-step
% situation C'' included in C' included in C with its functors Gamma',
% Gamma'' and Lambda = ST. Pages 52-53 name Gabriel.
%
% Page 54 is the wrapper of the third piece — "Théorèmes de descente de
% M. Artin – P. Cartier" — tallied "11 p", i.e. pages 56-66, so it runs past
% this batch and into batch 4. Pages 56-60 set out the gluing theorem: the
% functor f* = (i*, j*) : C -> C' x C'' under four conditions (a) to (d), the
% equivalence with the category of triples (G, H, u), the six-functor picture
% j*, j_*, j^!, i_!, i*, i_* of page 58, and Prop. 1 with its Corollary 1 —
% four conditions equivalent to the exactness of j* — proved on pages 59-60.
%
% The dating is the inventory's own, for the whole shelfmark; nothing on these
% pages carries a date.
%
% Pass: Opus 5 (claude-opus-5), 2026-08-22 — first pass, unchecked against the pages by a human.
\documentclass[11pt,a4paper]{article}
% Le préambule commun à toutes les transcriptions du fonds Grothendieck.
%
% Il tient en une idée : **l'incertitude se compose, elle ne se résout pas.**
% Une transcription qui lisse un mot douteux a détruit la seule chose pour
% laquelle on la faisait — un lecteur doit pouvoir savoir, à chaque endroit,
% ce qui était sur le feuillet et ce qui a été ajouté. Les six macros qui
% suivent sont donc l'appareil critique, et non de la décoration.
%
% Les mêmes commandes sont reconnues par `scripts/render.mjs`, qui produit la
% vue de lecture du site. Une macro ajoutée ici doit l'être aussi là-bas,
% faute de quoi le rendu échouera bruyamment — ce qui est voulu : un rendu
% silencieusement faux est pire qu'un rendu absent.

\ProvidesPackage{grothendieck}[2026/08/08 Transcriptions du fonds Grothendieck]

\usepackage{amsmath,amssymb,amsthm}
\usepackage{mathrsfs}
\usepackage{stmaryrd}
% Les diagrammes commutatifs sont partout dans ce fonds : c'est la langue
% même de la géométrie algébrique de Grothendieck, et une transcription qui
% les rendrait en prose ne transcrirait rien.
\usepackage{tikz-cd}
% Français par défaut — c'est la langue des feuillets — mais l'anglais est
% chargé aussi : la traduction et le résumé sont des documents anglais, et la
% typographie française (espaces avant les deux-points) y serait une faute.
% Ces documents basculent par \selectlanguage{english} en tête de corps.
\usepackage[english,french]{babel}
\usepackage{fontspec}
\usepackage{geometry}
\geometry{a4paper,margin=25mm,marginparwidth=28mm}
\usepackage{marginnote}
\usepackage{xcolor}
% ulem, not soul. `soul`'s \st and \ul are built on a character-by-character
% scan that XeLaTeX's font machinery breaks: the document fails with an
% undefined control sequence rather than degrading. ulem does the same two jobs
% and is Unicode-engine safe. `normalem` stops it hijacking \emph, which the
% transcriptions use for Grothendieck's own emphasis.
\usepackage[normalem]{ulem}
\usepackage{hyperref}
\hypersetup{colorlinks=true,linkcolor=black,urlcolor=black,pdfborder={0 0 0}}

% --- Métadonnées du lot -----------------------------------------------------
% Reprises telles quelles par le script de rendu, qui les affiche en tête de
% la vue de lecture. Elles fixent ce que le fichier prétend couvrir : sans
% elles, une transcription flotte sans point d'attache dans un fonds de seize
% mille feuillets.

\newcommand{\folder}[1]{\gdef\grothFolder{#1}}
\newcommand{\batch}[1]{\gdef\grothBatch{#1}}
\newcommand{\leaves}[2]{\gdef\grothFirst{#1}\gdef\grothLast{#2}}
\newcommand{\foldertitle}[1]{\gdef\grothFoldertitle{#1}}
\gdef\grothFolder{?}\gdef\grothBatch{?}\gdef\grothFirst{?}\gdef\grothLast{?}\gdef\grothFoldertitle{}

% --- Le résumé --------------------------------------------------------------
%
% Il ouvre la lecture modernisée, sous le titre « Résumé », et il est composé
% en petit. Ce n'est pas un ornement : c'est le seul passage du document qui
% ne soit pas des mathématiques, et un lecteur doit pouvoir le reconnaître
% avant de l'avoir lu — puis le sauter s'il connaît déjà le sujet, ou s'y
% arrêter s'il ne le connaît pas. Le filet qui le suit marque la reprise du
% corps.

\newenvironment{resume}
  {\small\setlength{\parskip}{0.45em}}
  {\par\medskip\hrule\medskip}

% --- Mots-clés ---------------------------------------------------------------
%
% En anglais, en clôture du résumé de la lecture modernisée : le vocabulaire
% moderne sous lequel chercher ce que les feuillets construisent. Le manifeste
% du site (`npm run manifest`) les extrait pour étiqueter la cote entière —
% cette ligne est la source unique des tags, il n'y a pas de fichier à côté.
\newcommand{\keywords}[1]{\par\smallskip\noindent
  {\footnotesize\textsc{Keywords} --- \textit{#1}}\par}

% --- L'appareil critique ----------------------------------------------------

% Le numéro de feuillet, en marge. C'est l'ancre qui permet de régler un
% désaccord sur une formule en regardant *un* feuillet plutôt qu'une chemise
% de six cents. Le site s'en sert aussi pour tourner les pages du fac-similé
% quand on fait défiler la transcription.
\newcommand{\leaf}[1]{%
  \marginnote{\footnotesize\color{gray!70}\textsf{#1}}%
  \label{leaf:#1}\ignorespaces}

% Illisible : on ne devine pas. Un mot inventé qui se lit comme les autres est
% le pire résultat possible.
\newcommand{\ill}{\textcolor{red!60!black}{[\,\ldots\,]}}

% Lecture douteuse : le mot est donné, et signalé comme incertain.
\newcommand{\uncertain}[1]{\uline{#1}}

% Ajout de l'éditeur — une lettre manquante, un mot sous-entendu.
\newcommand{\add}[1]{\textcolor{blue!50!black}{[#1]}}

% Biffé par Grothendieck lui-même. Ce qu'il a rayé fait partie du document :
% une rature dit souvent où la pensée a changé de direction.
\newcommand{\struck}[1]{\sout{#1}}

% Note du transcripteur, sur l'état matériel du feuillet.
\newcommand{\note}[1]{\par\smallskip{\small\color{gray!60!black}\textit{[#1]}}\par\smallskip}

% Note marginale de Grothendieck, distincte du corps du texte.
\newcommand{\marginal}[1]{\par\smallskip\noindent\hspace{1em}%
  {\small\color{gray!60!black}#1}\par\smallskip}

% --- Titre ------------------------------------------------------------------

\AtBeginDocument{%
  \begin{center}
    {\large\bfseries Fonds Alexandre Grothendieck --- cote n\textsuperscript{o}~\grothFolder}\\[2pt]
    {\itshape\grothFoldertitle}\\[4pt]
    {\small feuillets \grothFirst--\grothLast\ (lot \grothBatch)}\\[2pt]
    {\footnotesize\color{gray!60!black}Transcription établie d'après le fac-similé numérisé
     par l'Université de Montpellier.\\
     Les lectures incertaines sont soulignées, les illisibles notées [\,\ldots\,],
     les ajouts entre crochets.}
  \end{center}
  \vspace{1em}}

\endinput


\folder{19}
\batch{3}
\pages{41}{60}
\dating{[à partir de 1958-1973]}
\watermark{Édition de démonstration}
\foldertitle{Topos : notes manuscrites (s.d.)}

\begin{document}

\section*{Remarques — SGA 4 XII (pages 41 et 42)}

\page{41}
\note{feuille au crayon, très pâle ; la lecture reste lacunaire d'un bout à
l'autre}

XII

\uncertain{Dans} 6.5.

\emph{Remarque} (les hypothèses de quasi-compacité et quasi-séparation ne
servent que bien \struck{\ill{}} pour prouver l'équivalence de la condition
a) avec les conditions \ill{}, en supposant plus faible b), b\textsuperscript{$'$})
etc.

L'équivalence de ces conditions entre elles, et les implications de a)
\struck{\ill{}} dans b) \ill{} des 6.5.\ (i), \ill{} utilise \ill{} cette
hypothèse et mérite d'être \ill{} aux bonnes fins \ill{}. Dans 6.6.\ également,
on n'a \ill{} besoin de la quasi-compacité et quasi-séparation, \ill{} de plus,
la catégorie de tous les faisceaux \ill{}. Comme on l'a prouvé \ill{} d)
l'existence s'y \uncertain{complique} \ill{} une telle sous-catégorie —
manifestement on n'a pas d'ennui \ill{} tout abstrait. Dans le lemme 6.7.\ de
démonstration également, on n'utilise pas \ill{} quasi-compacité et
quasi-séparation, sauf, \struck{soi-disant}, bien pour : la limite \ill{} des
faisceaux en $\mathbf{Z}/n\mathbf{Z}$-Modules, — ou c'est sans doute inutile
\struck{\ill{}} — et on le \struck{verra}, de la même, en utilisant
\uncertain{XIV} 5.2., \ill{} déduit que l'\ill{} de changement de \ill{}
injectif — mais \ill{} directement.

\page{42}
\note{feuillet dactylographié de l'exposé ; deux traits diagonaux en biffent la
partie (iii), et l'auteur porte au bas de la page la note qui suit}

Adapter la démonstration de XII 8.3.\ (iii) ; ou se ramener si l'on veut au
cas $\uncertain{\pi_{0}}$ \emph{surjectif}, en remplaçant $\overline{X}$ par
$\overline{X} \amalg Y$ !

\section*{Foncteurs localisant et catégories de fractions (pages 43 à 53)}

\page{43}
Titre porté sur la chemise : « Foncteurs localisant et catégories de
fractions ».
\note{une main d'archiviste a pointé « 4 p + 5 p » en marge}

\page{44}
3) \quad
\begin{tikzcd}
A \arrow[r, bend left=20, "T"] & B \arrow[l, bend left=20, "S"]
\end{tikzcd}
\quad $\mathrm{Hom}(x, Sy) \simeq \mathrm{Hom}(Tx, y)$

$x \in \mathrm{Ob}\, A$, conditions équivalentes si $S$ pleinement fidèle :
\begin{enumerate}
\item[a)] $x$ est isomorphe à un objet de la forme $Sy$ ;
\item[b)] $x \xrightarrow{\ \sim\ } STx$ ;
\item[c)] $\forall\, x' \in A$, \quad
      $\mathrm{Hom}(x', x) \xrightarrow{\ \sim\ } \mathrm{Hom}(Tx', Tx)$.
\end{enumerate}

(b $\Rightarrow$ a sans hyp.\ sur $S$, a $\Rightarrow$ b si on suppose $S$
pleinement fidèle).

4) \quad
\begin{tikzcd}
A \arrow[r, bend left=20, "T"] & B \arrow[l, bend left=20, "S"]
\end{tikzcd}
\quad $\mathrm{Hom}(x, Sy) \simeq \mathrm{Hom}(Tx, y)$.

Conditions équivalentes :
\begin{enumerate}
\item[a)] $S$ pleinement fidèle ;
\item[b)] $TS \simeq \mathrm{id}_{B}$ isomorphisme ;
\item[c)] $T$ est un foncteur de passage au quotient, i.e.\ si on considère
      l'ens.\ $M$ des $m \in \mathrm{Fl}\, A$ tels que $Tm \in \mathrm{Isom}(B)$,
      et le foncteur $A/M \to B$, ceci est une équivalence de catégories.
\end{enumerate}

\page{45}
2) \quad $A \xrightarrow{\ T\ } A/C = B$

Conditions équivalentes sur $T$ i.e.\ sur $C$ :
\begin{enumerate}
\item[a)] $T$ admet un foncteur adjoint $S$ :
      $\mathrm{Hom}(Tx, y) \simeq \mathrm{Hom}(x, Sy)$ ;
\item[b)] $\forall\, x \in A$, il existe un plus grand sous-objet $x'$ de $x$
      qui est $\in C$, et si ce dernier est nul, il existe un monomorphisme de
      $x$ dans un objet $C$-fermé ;
\item[c)] $\forall\, x \in \mathrm{Ob}\, A$, il existe un hom.\
      $x \to \widetilde{x}$ qui est un $C$-isom., avec $\widetilde{x}$
      $C$-fermé.
\end{enumerate}

Alors le hom.\ en c) est \uncertain{unique} ; \ill{} puisque c'est aussi
$x \to STx = Lx$.

1) \quad $A \xrightarrow{\ T\ } A/C = B$, \quad $x \in \mathrm{Ob}\, A$.

Conditions équivalentes sur $x$ (pour $C$ i.e.\ $T$ donnée) :
\begin{enumerate}
\item[a)] $\forall\, x' \in A$, \quad
      $\mathrm{Hom}(x', x) \xrightarrow{\ \sim\ } \mathrm{Hom}(Tx', Tx)$ ;
\item[b)] $x' \to x''$ un $C$-isom.\ \ill{}, on a
      $\mathrm{Hom}(x'', x) \xrightarrow{\ \sim\ } \mathrm{Hom}(x', x)$ ;
\item[c)] tout sous-objet de $x$ qui est $\in C$ \ill{}, et \ill{}
      $x \subset x'$ tel que $x'/x \in C$ \ill{} $x$ \ill{} facteur direct.
\end{enumerate}
\marginal{\ill{}}
\note{une insertion de l'auteur, écrite en diagonale, se rattache à b) ; elle
est illisible}

\page{46}
$C$ catégorie abélienne avec $\varinjlim$ exactes et
\uncertain{générateurs}. $C'' \subset C'$ deux sous-catégories abéliennes
\emph{localisantes}, i.e.
\begin{enumerate}
\item[1\textsuperscript{o})] épaisses ;
\item[2\textsuperscript{o})] stables par $\varinjlim$.
\end{enumerate}

Foncteurs
\begin{align*}
  \Gamma'(F) &= \text{plus grand sous-objet de } F \in \mathrm{Ob}\, C' \\
  \Gamma''(F) &= \text{idem} \in \mathrm{Ob}\, C''
\end{align*}
\note{la seconde ligne est portée sur la page par un trait de rappel, et non
récrite}
\marginal{définir comme adjoints de $i'_{*}$, $i''_{*}$ ; noter $i'^{!}$,
$i''^{!}$}

\[
  \Gamma'/\Gamma''(F) = \Gamma'(F)/\Gamma''(F)
\]

Foncteur \quad
\begin{tikzcd}
C' \arrow[r, bend left=20, "T"] & C'/C'' \arrow[l, bend left=20, "S"]
\end{tikzcd}
\quad et foncteur section,

d'où foncteur $\Lambda = ST : C' \to C'$, avec
$\mathrm{id}_{C'} \to \Lambda$, \uncertain{caractérisé} par
\begin{enumerate}
\item[1\textsuperscript{o})] $x' \to \Lambda(x')$ est un $C''$-isom. ;
\item[2\textsuperscript{o})] $\Lambda(x')$ est « $C''$-fermé », i.e.\ pour tout
      hom.\ $y' \to z'$ dans $C'$,
      $\mathrm{Hom}_{C'/C''}(z', \Lambda x') \to
       \mathrm{Hom}_{C'/C''}(y', \Lambda x')$ est bijectif.
\end{enumerate}
\marginal{distinguer bien \ill{} $\Gamma'/\Gamma''$ \ill{} foncteur \ill{} $C$}

\textbf{Lemme 1.} $\Gamma'$, $\Gamma''$, $\Lambda$ existent ; \ill{},
$\Gamma'/\Gamma''$ hom.\ \ill{}
\[
  \Gamma'(F) \longrightarrow \Lambda\,\Gamma'(F)
\]
\ill{} Pour $\Lambda$, \ill{}. Pour $\Gamma'$ on note \ill{}
$\mathrm{Hom}_{C}\bigl( i'_{*}(x'), x \bigr) \simeq
 \mathrm{Hom}_{C'}\bigl( x', \Gamma'(x) \bigr)$, OK.

\page{47}
Soit $x \to y$ un mono, considérons $x''$, $y''$ les plus grands sous-objets
de $x$, $y \in \mathrm{Ob}\, C''$. Il faut prouver que

\begin{tikzcd}
x \arrow[r, hook] & y \\
x'' \arrow[r, hook] \arrow[u, hook] & y'' \arrow[u, hook]
\end{tikzcd}

$x'' = x \cap y''$, or c'est trivial.

\textbf{Cor.} $\Gamma'(x)/\Gamma''(x) \hookrightarrow
R^{0}(\Gamma'/\Gamma'')(x)$.

Considérons pour $x \in \mathrm{Ob}\, C$ variable
\[
  \Gamma'x \longrightarrow \Lambda\,\Gamma'x
\]
d'où un hom.\ $\Gamma'/\Gamma''(x) \longrightarrow \Lambda\,\Gamma'(x)$, et
comme $\Lambda\Gamma'$ est exact à gauche :

\begin{tikzcd}
\Gamma'x \arrow[r] \arrow[d] & \Lambda\Gamma' \\
R^{0}(\Gamma'/\Gamma'') \arrow[ur] &
\end{tikzcd}

\textbf{Prop.} C'est là un isom.\ (?)
\marginal{Faux en général ! OK si les injectifs de $C''$ sont
\uncertain{injectifs dans $C$}}

équivalent :

\textbf{Cor.} Pour tout $x \in \mathrm{Ob}\, C$ :
\begin{enumerate}
\item[a)] $R^{0}(\Gamma'/\Gamma'')(x)$ est $C''$-fermé [dans $C'$, ou dans
      $C$, c'est pareil \ill{}] ;
\item[b)] $\Gamma'x/\Gamma''(x) \to R^{0}(\Gamma'/\Gamma'')(x)$ est un
      $C''$-isom.
\end{enumerate}
\ill{}

\textbf{Cor.} Si $x$ est injectif, alors
$\Gamma'(x)/\Gamma''(x) \xrightarrow{\ \sim\ } \Lambda\,\Gamma'(x)$, i.e.\
$\Gamma'(x)/\Gamma''(x)$ est \struck{\ill{}} $C''$-fermé. En effet,
$x/\Gamma''(x)$ est $C''$-fermé, \struck{et} en effet, il est injectif par
loc.\ cit.\ cor.\ 2 et cor.\ 3.

\page{48}
Conditions à envisager :
\begin{enumerate}
\item[(i)] $C \xrightarrow{\ i_{*}\ } A$ transforme injectifs en injectifs ;
\item[(ii)] pour tout objet injectif $x$ de $A$, le plus grand sous-objet $x'$
      de $x$ qui est $\in \mathrm{Ob}\, C$ est encore injectif \emph{dans} $A$ ;
\item[(iii)] pour tout objet de $C$, une enveloppe injective dans $A$ est dans
      $C$ [i.e.\ $C$ stable par enveloppes injectives] — donc a.\ c.\
      injectifs dans $C$ ;
\item[(iv)] \struck{pour tout $x$} l'homomorphisme naturel
      \[
        R^{0}\bigl( \mathrm{id} / i_{*} i^{!} \bigr) \longrightarrow ST
      \]
      est un isomorphisme, i.e.\ pour $x$ objet injectif de $A$,
      \struck{si} si $x' = i_{*} i^{!}(x)$ le plus grand sous-objet
      $\in \mathrm{Ob}\, C$, alors $x/x'$ est $C''$-fermé
      \note{la page n'a introduit que $C$ ; le double prime est celui des
      pages précédentes} ;
\item[(v)] $T : A \to A/C$ transforme injectifs en injectifs.
\end{enumerate}
\marginal{cf.\ Gabriel, cor.\ 3}

(i) $\Longrightarrow$ (ii) : car \struck{\ill{}} pour tout $x$,
$x' = i_{*} i^{!}(x)$, or $i^{!}(x)$ [étant adjoint de $i_{*}$, \ill{}
conditions \ill{}] transforme injectifs en injectifs.

(ii) $\Longrightarrow$ (iii) : car si $x \in C$, soit $\overline{x}$ une
enveloppe injective de $x$ dans $A$, et considérons
$i_{*}\bigl( i^{!}(\overline{x}) \bigr)$, c'est un objet injectif de $C$
contenant $x$, et une extension essentielle, donc [car \ill{}] donc c'est une
enveloppe injective de $x$ dans $C$. Si $\overline{x}$ est injectif, il
\ill{} égal à $i_{*} i^{!}(\overline{x})$. Si (ii) est vérifié,
$i_{*} i^{!}(\overline{x})$ est injectif, d.\ c.\ $i(x)$ injectif.

\page{50}
\textbf{Lemme.} Soit $x \in \mathrm{Ob}\, C$, et soit $\overline{x}$ une
enveloppe injective dans $A$, alors $i^{!}(\overline{x})$ est une enveloppe
injective de $x$ dans $C$, [donc en particulier \ill{} $x$ est \ill{} injectif
dans $C$].

\textbf{Cor.} Tout objet injectif de $C$, aux isom.\ près,
\struck{\ill{}} est de la forme $i^{!}(y)$, $y$ un objet injectif de $A$.

\textbf{Cor.} (i) $\Longleftrightarrow$ (ii) $\Longleftrightarrow$ (iii).

(ii) $\Longrightarrow$ (iv) : car si $x$ est injectif, on a \ill{}, car
$i_{*} i^{!}(x)$ dans $x$ est \struck{injectif} $C$-fermé par loc.\ cit.\
cor.\ 2.

\page{52}
$\mathrm{Hom}(x, Sy) \simeq \mathrm{Hom}(Tx, y)$

\textbf{Prop.} \quad
\begin{tikzcd}
A \arrow[r, bend left=20, "T"] & B \arrow[l, bend left=20, dashed, "S"]
\end{tikzcd}
\quad $T$ foncteur « passage au quotient » pour une catégorie de fractions.

Soit $x \in \mathrm{Ob}\, A$, $y = T(x)$. Propriétés équivalentes :
\begin{enumerate}
\item[a)] $x$ est libre au dessus de $y$, i.e.\ $\forall\,
      x' \in \mathrm{Ob}\, A$,
      \[
        \mathrm{Hom}_{A}(x', x) \xrightarrow{\ \sim\ }
        \mathrm{Hom}_{B}(Tx', Tx) ;
      \]
\item[b)] si $x' \to x''$ dans $A$ est un « $C$-isomorphisme » (i.e.\
      $Tx' \to Tx''$ un isomorphisme) alors
      $\mathrm{Hom}_{A}(x'', x) \longrightarrow \mathrm{Hom}_{A}(x', x)$
      est une bijection ;
\item[c)] [si $B = A/C$, $A$ catégorie abélienne, $C$ sous-catégorie abélienne
      épaisse] tout sous-objet de $x$ « \ill{} » : $\in \mathrm{Ob}\, C$ est
      nul, et $x$ est facteur direct dans tout $y \supset x$ tel que
      $y/x \in \mathrm{Ob}\, C$ ;
\item[d)] [si l'adjoint $S$ de $T$ existe] $x \xrightarrow{\ \sim\ } STx$ est
      un isomorphisme ;
\item[e)] \struck{[si l'adjoint $S$ de $T$ existe]} $x$ est dans l'image
      essentielle de $S$, i.e.\ isomorphe à un objet de la forme $Tz$
      [$z \in \mathrm{Ob}\, C$]
      \note{la forme attendue serait plutôt $Sz$ ; la page porte $Tz$}.
\end{enumerate}

\page{53}
\struck{\textbf{Corollaire.} Soit $x \in \mathrm{Ob}\, A$, \struck{\ill{}} et
$x \to x'$ une flèche de $A$. Conditions équivalentes : a) $Tx
\xrightarrow{\ \sim\ } Tx'$, et $x'$ est libre sur $Tx'$. Cf.\ prop.\ préc.}
\note{la proposition entière est encadrée puis biffée en croix}

\textbf{Corollaire 1.} Pour que le foncteur adjoint $S$ existe, il f.\ et s.\
que pour tout $x \in \mathrm{Ob}\, A$ existe une flèche $x \to x'$ telle que
\begin{enumerate}
\item[(i)] $Tx \to Tx'$ soit un isom. ;
\item[(ii)] $x'$ est libre sur $Tx'$, i.e.\ satisfait aux conditions de la
      prop.
\end{enumerate}

D'ailleurs, Gabriel prouve

\textbf{Corollaire 2.} \struck{Si $C$} Dans le cas $A$ abélienne, $B = A/C$
($C$ sous-catégorie épaisse), pour que \struck{$ST$} $S$ existe, il f.\ et s.\
que
\begin{enumerate}
\item[a)] $\forall\, x \in \mathrm{Ob}\, C$, existe un plus grand sous-objet
      $x''$ de $x$ qui soit $\in \mathrm{Ob}\, C$
      \note{la page porte $\mathrm{Ob}\, C$ ; c'est $\mathrm{Ob}\, A$ qui est
      en jeu} ;
\item[b)] [si $x'' = 0$,] il existe un monomorphisme de $x$ \struck{\ill{}}
      dans un objet « $C$-fermé ».
\end{enumerate}
\marginal{Il \ill{} la première de ces conditions \ill{} conditions \ill{}
l'existence de $STx$, $x \to \widetilde{x}$}

Si $C$ est \uncertain{à} $\varinjlim$ \uncertain{exactes} et \ill{}
d'intersections, la condition a) est \ill{}, et signifie \ill{} que $C$ est
stable par $\varinjlim$. Alors $T$ transforme injectifs en injectifs.

\section*{Théorèmes de descente de M. Artin – P. Cartier (pages 54 à 60)}

\page{54}
Titre porté sur la chemise : « Théorèmes de descente de M. Artin – Cartier ».
\note{une main d'archiviste a pointé « 11 p » en marge, ce qui est exactement
le compte des pages 56 à 66 ; le texte déborde donc ce lot}

\page{56}
\note{l'auteur numérote lui-même cette page 1, et les suivantes 2, 3, 4}

\emph{Th.\ de M.\ Artin et P.\ Cartier}

Soient $C$, $C'$, $C''$ des catégories,
\[
  f^{*} = (i^{*}, j^{*}) : C \longrightarrow C' \times C'' = D
\]
un foncteur. On suppose que $i^{*}$, $j^{*}$ ont des adjoints $i_{*}$,
$j_{*}$ :
\[
  \mathrm{Hom}(F, i_{*}G) = \mathrm{Hom}(i^{*}F, G), \qquad
  \mathrm{Hom}(F, j_{*}H) = \mathrm{Hom}(j^{*}F, H)
\]
donc que $f^{*}$ a un adjoint $f_{*}$
\[
  f_{*}(G, H) \simeq i_{*}(G) \times j_{*}(H)
\]
\begin{align*}
  \bigl[\ \mathrm{Hom}\bigl( F, f_{*}(G,H) \bigr)
    &\simeq \mathrm{Hom}\bigl( F, i_{*}(G) \bigr) \times
            \mathrm{Hom}(F, j_{*}H) \\
    &\simeq \mathrm{Hom}(i^{*}F, G) \times \mathrm{Hom}(j^{*}F, H) \\
    &\simeq \mathrm{Hom}\bigl( f^{*}F, (G,H) \bigr)\ \bigr]
\end{align*}

On suppose de plus
\begin{enumerate}
\item[(a)] $i^{*} i_{*} G \simeq G$ i.e.\ $i_{*}$ est pleinement fidèle ;
      $j^{*} j_{*} H \simeq H$ i.e.\ $j_{*}$ est pleinement fidèle ;
\item[(b)] $i^{*} j_{*}(H)$ est objet final de $C'$ pour tout
      $H \in \mathrm{Ob}\, C''$.

      [Alors $f^{*} f_{*}(G, H) \simeq \bigl( G, \varphi(G) \times H \bigr)$,
      et \struck{\ill{}} les « données de recollement » sur $(G, H)$ \ill{}
      $f^{*}$ \ill{} dans \ill{}
      $\bigl( G, \varphi(G) \times H \bigr) \longleftarrow (G, H)$ de la forme
      $\bigl( \mathrm{id}_{G}, (u, \mathrm{id}_{H}) \bigr)$, i.e.\
      $u : H \to \varphi(G)$ est arbitraire] ;
\item[(c)] les noyaux et co-noyaux existent dans $C$, et $f^{*}$ commute aux
      noyaux et co-noyaux [i.e.\ $i^{*}$, $j^{*}$ y commutent].
\end{enumerate}
\marginal{N.B.\ on pose $\varphi = j^{*} i_{*}$ \ill{}}

\page{57}
\begin{enumerate}
\item[(d)] $f^{*}$ est conservatif, i.e.\ si $\lambda \in \mathrm{Fl}\, C$ est
      tel que $i^{*}\lambda$ et $j^{*}\lambda$ soient des isom., alors
      $\lambda$ est un isom.
\end{enumerate}

Sous ces conditions, le foncteur naturel
\[
  C \xrightarrow{\ \theta\ } K
\]
[où $K$ est la \struck{\ill{}} catégorie des triples $(G, H, u)$,
$G \in \mathrm{Ob}\, C'$, $H \in \mathrm{Ob}\, C''$,
$u \in \mathrm{Hom}\bigl( H, \varphi(G) \bigr)$] est une \emph{équivalence de
catégories}.

\emph{Ex.} \struck{\ill{}} $C = $ faisceaux sur $X$ ; $C' = $ faisceaux sur
$U \subset X$ ($U$ ouvert) ; $C'' = $ faisceaux sur $Y \subset X$
($Y = X - U$).

\uncertain{équivalent} \ill{} \struck{\ill{}}
\begin{enumerate}
\item[d$'$)] pour tout $F \in \mathrm{Ob}\, C$, le diagramme
\end{enumerate}

\begin{tikzcd}
F \arrow[r] \arrow[d] & i_{*} i^{*} F \arrow[d] \\
j_{*} j^{*} F \arrow[r] & j_{*} j^{*} i_{*} i^{*} F
\end{tikzcd}

(i.e.

\begin{tikzcd}
F \arrow[r] \arrow[d] & i_{*}(G) \arrow[d] \\
j_{*}(H) \arrow[r] & j_{*}\bigl( \varphi(G) \bigr)
\end{tikzcd}

) est cartésien.

\page{58}
La situation

\begin{tikzcd}[column sep=huge]
C'' \arrow[r, "j_{*}" description] &
C \arrow[l, bend right=45, "j^{*}"'] \arrow[l, bend left=45, "j^{!}"]
  \arrow[r, "i^{*}" description] &
C' \arrow[l, bend right=45, "i_{!}"'] \arrow[l, bend left=45, "i_{*}"]
\end{tikzcd}

\note{les trois objets sont notés $H$, $F$, $G$ au-dessous de $C''$, $C$, $C'$}
(catégories abéliennes, foncteurs additifs)

\begin{align*}
  \mathrm{Hom}(F, j_{*}H) &= \mathrm{Hom}(j^{*}F, H) \\
  \mathrm{Hom}(j_{*}H, F) &\simeq \mathrm{Hom}(H, j^{!}F) \\
  \mathrm{Hom}(F, i_{*}G) &\simeq \mathrm{Hom}(i^{*}F, G) \\
  \mathrm{Hom}(i_{!}G, F) &\simeq \mathrm{Hom}(G, i^{*}F)
\end{align*}

$j_{!} = j_{*}$, $j^{*}$, $i^{*} = i^{!}$, $i_{!}$ exacts ; $j^{!}$, $i_{*}$
exacts à gauche.

$\varphi = j^{*} i_{*} : C'' \longleftarrow C'$, \quad $\varphi$ exact à
gauche.

\begin{align*}
  j_{*} j^{!} F &= \Gamma_{Y} F & j_{*} j^{*} F &= F^{Y} \\
  i_{*} i^{*} F &= \Gamma_{U} F & i_{!} i^{*} F &= F^{U}
\end{align*}

\[
  i^{*} j_{*} = 0, \qquad j^{!} i_{*} = 0, \qquad
  j^{*} i_{!} = 0, \qquad j^{!} i_{!} = 0
\]
\note{les trois premières portent une croix en marge}

\[
  j^{*} \longrightarrow \varphi\, i^{*}
\]

\textbf{Prop.\ 1.} Partons de $C'' \xrightarrow{\ j_{*}\ } C
\xrightarrow{\ i^{*}\ } C'$, où $j_{*}$ est pleinement fidèle, $i^{*}$ est un
passage au quotient exact [donc $i^{*} j_{*} = 0$], et
$j_{*} \simeq \mathrm{Ker}\, i^{*}$. Alors
\begin{enumerate}
\item[(i)] si $i_{!}$ existe [i.e.\ pl.\ fid.], alors $j^{*}$ existe, et on a
      la suite ex.
      \[
        i_{!} i^{*} F \longrightarrow F \longrightarrow j_{*} j^{*} F
        \longrightarrow 0
      \]
      De plus \quad
      $j^{*} j_{*} \simeq \mathrm{id}_{C''}$, \quad
      $i^{*} i_{!} \xrightarrow{\ \sim\ } \mathrm{id}_{C'}$, \quad
      $j^{*} i_{!} \simeq 0_{C''}$ ;
\item[(i bis)] si $i_{*}$ existe, alors $j^{!}$ existe, et on a
      \[
        i_{*} i^{*} F \longleftarrow F \longleftarrow j_{*} j^{!} F
        \longleftarrow 0
      \]
      De plus \quad
      $j^{!} j_{*} \simeq \mathrm{id}_{C''}$, \quad
      $i^{*} i_{*} \xrightarrow{\ \sim\ } \mathrm{id}_{C'}$, \quad
      $j^{!} i_{*} \simeq 0_{C'}$.
\end{enumerate}
\marginal{Question : si $j^{*}$ existe, est-il vrai que $i_{!}$ existe ?}

(N.B.\ En général, on ne peut \uncertain{y} ajouter les zéros \ldots)

\page{59}
L'énoncé (i bis) est dual à (i).

\textbf{Corollaire 1.} [Sous les cond.\ de prop.\ 1] Supposons que $i_{!}$
existe (condition (i) de la proposition 1). Les conditions suivantes sont
équivalentes :
\begin{enumerate}
\item[a)] $j^{*}$ exact, $i_{!} \simeq \mathrm{Ker}\, j^{*}$ [a fortiori,
      $i_{!}$ est \struck{\ill{}} exact] ;
\item[b)] $i_{!}$ a une image essentielle épaisse [épaisse à g.\ suffit] ;
\item[c)] $i_{!} i^{*}(F) \longrightarrow F$ est injectif pour tout $F$,
      i.e.\ la suite
      \[
        0 \longrightarrow i_{!} i^{*} F \longrightarrow F \longrightarrow
        j_{*} j^{*} F \longrightarrow 0
      \]
      est exacte ;
      \note{les deux termes extrêmes sont annotés « $F^{U}$ » et « $F^{Y}$ »}
\item[d)] $j^{*}$ est exact.
\end{enumerate}
\struck{Alors $i_{!} \simeq \mathrm{Ker}\, j^{*}$, en particulier $i_{!}$ est
exact.}

a) $\Longrightarrow$ b) trivial.

b) $\Longrightarrow$ c) : car $0 \to N \to i_{!} i^{*} F \to F$
$\Longrightarrow$ $0 \to i^{*}N \to i^{*} i_{!} i^{*} F
\xrightarrow{\ \sim\ } i^{*}F$ [$i^{*}$ exact] donc \ill{}, donc
$i^{*}(N) = 0$ \struck{\ill{}}, or $N = i_{!}(M)$ par b),
$i^{*} i_{!}(M) = 0 \Rightarrow M = 0 \Rightarrow N = 0$ [$i^{*}$ passage au
quotient].

c) $\Longrightarrow$ a) : \struck{\ill{}} supposons
$0 \to F \to F' \to F'' \to 0$ \ill{}

\struck{a) $\Longrightarrow$ d) trivial}

d) $\Longrightarrow$ b) :

\begin{tikzcd}[column sep=small, row sep=small, nodes={font=\scriptsize}]
0 \arrow[r] & i_{!} i^{*} F \arrow[r] \arrow[d] & F \arrow[r] \arrow[d] &
  j_{*} j^{*} F \arrow[r] \arrow[d] & 0 \\
0 \arrow[r] & i_{!} i^{*} F' \arrow[r] \arrow[d] & F' \arrow[r] \arrow[d] &
  j_{*} j^{*} F' \arrow[r] & 0 \\
0 \arrow[r] & i_{!} i^{*} F'' \arrow[r] & F'' & &
\end{tikzcd}

en utilisant l'injectivité de $i_{!} i^{*} F'' \to F''$, on \uncertain{trouve}
$j_{*} j^{*} F \to j_{*} j^{*} F'$ est injectif, donc $j^{*}F \to j^{*}F'$
aussi, et $j^{*}$ ex.\ (car il est déjà exact à droite).

Soit $F = i_{!}G$, alors $i_{!} i^{*} F \longrightarrow F$ est un isom.\
[$i^{*}$ passage au quotient] donc $j_{*} j^{*} F = 0$ \struck{\ill{}} par la
suite exacte (i), \struck{donc $j^{*}F = 0$.} Réciproquement, si
$j^{*}F = 0$, alors par (i), on trouve
$i_{!} i^{*} F \xrightarrow{\ \sim\ } F$, i.e.\ $F \in$ \struck{\ill{}}
$\mathrm{Im}\, i_{!}$.
\marginal{\ill{} \uncertain{prop.\ 1}}

\page{60}
\struck{a) $\Longrightarrow$ d) trivial}

d) $\Longrightarrow$ c) : car \struck{\ill{}} $0 \to N \to i_{!} i^{*} F \to
F$ \ill{} $\Longrightarrow$ $j^{*}N \to j^{*} i_{!} i^{*} F$ \struck{\ill{}}
[par d)], $= 0$, d'où $j^{*}N = 0$, or \ill{} $i^{*}N = 0$
$\bigl[ 0 \to i^{*}N \to i^{*} i_{!} i^{*} F \xrightarrow{\ \sim\ } i^{*}F
\bigr]$ donc $N = j_{*}M$, $M \in \mathrm{Ob}\, C''$, et comme
$H = j^{*} j_{*} H \simeq j^{*}N = 0$, on aura $N = 0$, cqfd.

\struck{Cor.\ 2}

\textbf{Corollaire 1 bis.} dual. (Les conditions ci-dessus ne sont pas
vérifiées dans les situations en vue, en particulier $j^{!}$ ne sera pas
exact, $\mathrm{Im}\, i_{*}$ ne sera pas épaisse à droite,
$F \to i_{*} i^{*} F$ ne sera pas en général surjectif].

\textbf{Remarque 1.} \struck{\ill{}} Les conditions du cor.\ 1 ne résultent
pas du fait que $i_{!}$ est exact, comme on voit dans la situation décrite par
la localisation des modules \ldots

\textbf{Remarque 2.} Sous les conditions équiv.\ du cor.\ 1, \ill{} la
situation

\begin{tikzcd}[column sep=huge]
C'' \arrow[r, "j_{*}" description] &
C \arrow[l, bend right=35, "j^{*}"'] \arrow[r, "i^{*}" description] &
C' \arrow[l, bend right=35, "i_{!}"']
\end{tikzcd}

\ill{} \struck{\ill{}} \uncertain{symétrique} en $C'$, $C''^{\circ}$ :
\begin{enumerate}
\item[] $i^{*}$ passage au \struck{\ill{}} exact de noyau
      \uncertain{$i_{!}$}
      \note{la symétrie annoncée demanderait plutôt $j_{*}$ ici} ;
\item[] $j^{*}$ passage au quotient exact de noyau $i_{!}$ ;
\item[] $j_{*}$ et $j^{*}$ adjoints, \quad $i^{*}$ et $i_{!}$ adjoints.
\end{enumerate}

De façon précise, la situation

\begin{tikzcd}[column sep=huge]
C'^{\circ} \arrow[r, "i_{!}^{\circ}" description] &
C^{\circ} \arrow[l, bend right=35, "i^{*\circ}"'] \arrow[r, "j^{*\circ}" description] &
C''^{\circ} \arrow[l, bend right=35, "j_{*}^{\circ}"']
\end{tikzcd}

satisfait \struck{\ill{}} aux hyp.\ pour la situation au départ. Cette
symétrie va être détruite dans ce qui suit.

\end{document}
